\SourceLeaf{341}%
% A544-P0341-D01
\begin{list}{}{\setlength{\labelwidth}{9em}\setlength{\leftmargin}{9.5em}}
\item[$\alpha=-(\Upsilon\overline{\overline\lambda}+1)=-\Upsilon l$] für alle auf dieser Grenze, die in entgegengesetztem Sinne laufen
\item[$\alpha=\Upsilon l$] für alle Grenzligaturen innerhalb $\Omega'$, deren Richtung auf $0$ zugeht
\item[$\alpha=-\Upsilon l+1$] für alle Grenzligaturen innerhalb $\Omega'$, deren Richtung von $0$ abwärts geht.
\end{list}

% A544-P0341-D02
\Article{12}
Wir können nun die sämmtlichen vorkommenden $l$ (nach der letzten Manier) zu zweien combiniren, nemlich $l$ mit $\overline{\overline{\frac12m-l}}$, welche wir verbundene Ligaturen nennen wollen; eine einzige ist hiervon ausgenommen, welche isolirt steht oder mit ihrer verbundenen Ligatur identisch ist, nemlich diejenige, welche von
\[
\tfrac12(M-1)\pdot\frac{m}{2M}\quad\text{nach}\quad\tfrac12(M+1)\frac{m}{2M}\quad\text{läuft}
\]
für verbundene Ligaturen ist das $\alpha$ immer einerlei.
\BookRule{.3\textwidth}

% A544-P0341-D03
\par\bigskip\Needspace{7\baselineskip}
\pdfbookmark[1]{Theorie der biquadratischen Reste. III}{nachlass-biquadratic-III}
\begin{center}[III.]\par\medskip\textsc{Theorie der biquadratischen Reste.}\end{center}
\BookRule{.08\textwidth}
\renewcommand{\Article}[1]{\par\medskip\Needspace{4\baselineskip}\pdfbookmark[2]{Art. #1}{nachlass-biquadratic-III-art-#1}\centerline{\textbf{#1.}}\nobreak\smallskip}
\Article{1}
Kleinste Reste des Modulus $m=a+bi$ heissen die ganzen Zahlen $\mu=\alpha+\beta i$, für welche $\frac{\mu}{m}=x+yi$ so beschaffen ist, dass $x$ und $y$ positiv und kleiner als $1$ sind. Es kommt noch dazu der Rest $0$\begingroup\renewcommand{\thefootnote}{*)}\footnote{und wenn $a$ und $b$ etwa den gemeinschaftlichen Divisor $e$ haben, die Zahlen $\frac{m}{e},\frac{2m}{e},\frac{3m}{e}\dotstwo\frac{(e-1)m}{e}$. Jedoch wollen wir diesen Fall vorerst von der Untersuchung ausschliessen.}\endgroup. Ihre Anzahl ist $=aa+bb$.

% A544-P0341-D04
\Article{2}
In sofern $aa+bb$ ungerade ist, wird $aa+bb$ von der Form $4n+1$ sein. Den kleinsten Rest $0$ ausgeschlossen, theilen sich die übrigen in vier Classen. Zur ersten Classe $f$ zählen wir diejenigen, wo $x$ und $y$ kleiner als $\frac12$ sind,
\SourceLeaf{342}%
% A544-P0342-D01
\[
\begin{array}{lcll}
\text{die zweite}&f'&\text{wo }x>\tfrac12,&y<\tfrac12\\
\text{dritte}&f''&x>\tfrac12,&y>\tfrac12\\
\text{vierte}&f'''&x<\tfrac12,&y>\tfrac12
\end{array}
\]
Man erhält alle Reste
\[
\begin{array}{rcl}
f'&\text{aus}&if+m\\
f''&\text{aus}&-f+(1+i)m\\
f'''&\text{aus}&-if+im
\end{array}
\]

% A544-P0342-D02
\Article{3}
Es sei $M$ eine andere Zahl, die mit $m$ keinen Factor gemein hat, so wird
\[
M^{aa+bb-1}\equiv1\ (\mathrm{mod.}\ m)
\]
sein: folglich $M^{\frac14(aa+bb-1)}$ entweder $\equiv1$, oder $\equiv i$, oder $\equiv-1$, oder $\equiv-i$ d.\,i. $\equiv i^{\varepsilon}$, wo $\varepsilon$ eine der vier Zahlen $0,1,2,3$ vorstellt. Im ersten Fall wird $M$ biquadratischer Rest von $m$ sein, mithin auch quadratischer. Im dritten ist $M$ quadratischer aber nicht biquadratischer Rest; im zweiten und vierten sowohl quadratischer als biquadratischer Nichtrest. Wir nennen dies $\varepsilon$, wovon die biquadratische Modalität der Zahl $M$ in Beziehung auf den Modulus $m$ abhängt, den Decidenten von $M$ beim Modulus $m$. Die Induction lehrt folgenden schönen Lehrsatz. „Sind $M$ und $m$ ungerade Primzahlen von der Form $1+(2+2i)\mu$, so dass $\mu$ eine ganze Zahl ist, so ist die Differenz der beiden Decidenten, von $M$ beim Modulus $m$, und von $m$ beim Modulus $M$ entweder $=0$ oder $=2$; das erstere, wenn wenigstens eine der Zahlen $m,M$ von der Form $1+4N$ ist; das andere, wenn beide von der Form $1+2i+4N$ sind.“ Dies Theorem der Reciprocität ist dem bei den Quadratischen Resten bei blos reellen Zahlen analog.

% A544-P0342-D03
\Article{4}
Man multiplicire alle Zahlen $f$ mit $M$, und suche deren kleinste Reste nach dem Modulus $m$. Es seien darunter
\[
\begin{array}{cll}
\alpha&\text{zu }f&\text{gehörig}\\
\beta&f'&\\
\gamma&f''&\\
\delta&f'''&
\end{array}
\]
so ist $\varepsilon\equiv\beta+2\gamma+3\delta\ (\mathrm{mod.}\ 4)$.

\SourceLeaf{343}%
% A544-P0343-D01
Beweis. Der Inbegriff derjenigen Zahlen aus $f$, deren Producte mit $M$ Reste zu $f$ gehörig geben, sei $g$; der Inbegriff derjenigen, deren Producte Reste aus $f'$ geben, sei $g'$, und ebenso $g''\pdot g'''$; so werden die kleinsten Reste von
\[
-ig'M,\quad -g''M,\quad ig'''M
\]
alle in $f$ enthalten, und sowohl unter sich als von den kleinsten Resten der Producte $gM$ verschieden sein, folglich das Product aus allen
\[
gM,\quad -ig'M,\quad -g''M,\quad +ig'''M
\]
dem Producte aller $f$ congruent sein, mithin auch dem Producte aller $g,g',g'',g'''$. Jenes Product ist aber gleich dem Producte aus allen $g,g',g'',g'''$ in
\[
M^{\alpha}\pdot(-iM)^{\beta}\pdot(-M)^{\gamma}\pdot(iM)^{\delta}
\]
also dies letzte Product $\equiv1$
\[
\begin{aligned}
\text{folglich}\qquad&M^{\alpha+\beta+\gamma+\delta}(-i)^{\beta}(-1)^{\gamma}i^{\delta}\equiv1\\
\text{oder}\qquad&M^{\alpha+\beta+\gamma+\delta}\equiv i^{\beta}(-1)^{\gamma}(-i)^{\delta}\equiv i^{\beta+2\gamma+3\delta}
\end{aligned}
\]
woraus der Lehrsatz von selbst folgt.

% A544-P0343-D02
\Article{5}
Die Entscheidung, ob der kleinste Rest einer Zahl $N$ nach dem Modulus $m$ zur Classe $f,f',f''$ oder $f'''$ gehöre, ist leicht. Ist nemlich $\omega$ die in $\frac{N}{m}$ enthaltene ganze Zahl, so wird jener Rest $=N-\omega m$ sein, und also zu $f,f',f''$ gehören, jenachdem
\[
\frac{N}{m}-\omega=x+iy
\]
gesetzt
\[
\begin{array}{ll}
x<\tfrac12,&y<\tfrac12\\
x>\tfrac12,&y<\tfrac12\\
x>\tfrac12,&y>\tfrac12\\
x<\tfrac12,&y>\tfrac12
\end{array}
\]
ist. In diesen $4$ Fällen wird der Reihe nach die in $\frac{2N}{m}$ enthaltene ganze Zahl folgende sein
\SourceLeaf{344}%
% A544-P0344-D01
\[
\begin{array}{l}
2\omega\\
2\omega+1\\
2\omega+1+i\\
2\omega+i
\end{array}
\]
Hieraus ist klar, dass der kleinste Rest von $N$ nach dem Modulus $m$ zu $f,f',f'',f'''$ gehören werde, je nachdem die in $\frac{2N}{m}$ enthaltene ganze Zahl $=\xi+\eta i$ gesetzt
\[
\begin{array}{ll}
\xi\ \text{gerade}&\eta\ \text{gerade}\\
\xi\ \text{ungerade}&\eta\ \text{gerade}\\
\xi\ \text{ungerade}&\eta\ \text{ungerade}\\
\xi\ \text{gerade}&\eta\ \text{ungerade}
\end{array}
\]

% A544-P0344-D02
\Article{6}
Hienach findet sich der Decident von $M$ nach dem Modulus $m$ auf folgende Art. Man suche die ganzen Zahlen, die in allen einzelnen $\frac{2fM}{m}$ enthalten sind. Diese allgemein durch $x+yi$ bezeichnet, lasse man ganz aus der Acht, diejenigen, wo $x$ und $y$ beide gerade sind, rechne für jede derjenigen, wo $x$ ungerade und $y$ gerade ist, eins, entnehme für jede derjenigen, wo $x$ und $y$ beide ungerade sind, zwei, und drei für jede von denen, wo $x$ gerade, $y$ ungerade ist. Von der Summe aller dieser Zahlen nehme man den kleinsten Rest nach $4$, welcher der verlangte Decident sein wird. Wir drücken dies so aus
\[
\mathrm{Dec.}\frac{M}{m}=\Sigma n
\]
wo $[\frac{2fM}{m}]=x+yi$, $n=0$ zu setzen ist, wenn $x$ gerade $y$ gerade
\[
\begin{array}{cll}
1&x\ \text{ungerade}&y\ \text{gerade}\\
2&x\ \text{ungerade}&y\ \text{ungerade}\\
3&x\ \text{gerade}&y\ \text{ungerade}
\end{array}
\]
Kürze halber wollen wir $n$ durch die Characteristik $\Theta$ bezeichnen, $n=\Theta\frac{2fM}{m}$\begingroup\renewcommand{\thefootnote}{*)}\footnote{Um zu entscheiden, in welche Classe $M$ in Beziehung auf $m$ gehört, wählt man diejenigen Werthe von $k$ (unter den Zahlen $1,2,3\ldots p-1$) aus, wodurch $[\frac{2km'M}{p}]$ gerade wird, und addirt $-\Sigma[\frac{2km'}{p}]^2$ Nimmt man $k$ nur bis $\frac12p$, so hat man zu summiren
\[
-\Sigma\{[\frac{2km'M}{p}]^2+[\frac{2km'}{p}]^2\}
\]
für diejenigen Werthe von $[\frac{2km'M}{p}]$ die durch $1+i$ theilbar sind.}\endgroup.

\SourceLeaf{345}%
% A544-P0345-D01
\par\medskip\Needspace{4\baselineskip}\pdfbookmark[2]{Art. 7. 8.}{nachlass-biquadratic-III-art-7-8}\centerline{\textbf{7. 8.}}\nobreak\smallskip
Diese Regel ist allgemein, was für eine Zahl auch $M$ bedeute. Für den Fall, der zunächst den Gegenstand unserer Untersuchung ausmachen soll, wo $M$ ungerade und von der Form $1+(2+2i)N$ vorausgesetzt wird, ist eine etwas abgeänderte Vorschrift zweckmässiger.

% A544-P0345-D02
Man denke sich die Zahlen $f$ wiederum in $4$ Classen zerlegt; in die erste setzt man die $(h)$, deren Doppeltes sich auch noch in $f$ findet; in die zweite $h'$ zählen wir die, deren Doppelte $2h'$ zu $f'$ gehören, und ebenso $h''$ und $h'''$ bedeuten diejenigen, deren Doppelte zu $f''$ und $f'''$ gehören. Es ist also der Decident $\varepsilon$
\[
\varepsilon=\Sigma\Theta\frac{2hM}{m}+\Sigma\Theta\frac{2h'M}{m}+\Sigma\Theta\frac{2h''M}{m}+\Sigma\Theta\frac{2h'''M}{m}
\]
Den Complexus aller $2h$ und $-2h''+(1+i)m$ nennen wir $H$

den von allen $-i(2h'-m)$ und $i(2h'''-im)$ nennen wir $H'$

% A544-P0345-D03
$H$ und $H'$ umfassen also alle $f$, jene sind die geraden, diese die ungeraden.

Ferner sind folgende Relationen in Anwendung zu bringen
\[
\begin{aligned}
\Theta iN&=1+\Theta N\\
\Theta(-N)&=2+\Theta N\\
\Theta(-iN)&=3+\Theta N\\
\Theta(N+1)&=1-\Theta N\\
\Theta(N+1+i)&=2+\Theta N\\
\Theta(N+i)&=3-\Theta N
\end{aligned}
\]
folglich
\[
\begin{aligned}
\Theta\frac{(-2h'i+mi)M}{m}&=3-\Theta\frac{-2h'iM}{m}=-\Theta\frac{2h'M}{m}\\
\Theta\frac{(-2h''+m(1+i))M}{m}&=2+\Theta\frac{-2h''M}{m}=\Theta\frac{2h''M}{m}\\
\Theta\frac{(2h'''i+m)M}{m}&=1-\Theta\frac{2h'''iM}{m}=-\Theta\frac{2h'''M}{m}
\end{aligned}
\]
und
\[
\begin{aligned}
\varepsilon={}&\Sigma\Theta\frac{2hM}{m}-\Theta\frac{(-2h'i+mi)M}{m}\\
&+\Theta\frac{(-2h''+m(1+i))M}{m}-\Theta\frac{(2h'''i+m)M}{m}\\
={}&\Sigma\Theta\frac{HM}{m}-\Sigma\Theta\frac{H'M}{m}
\end{aligned}
\]

